% !TEX program = xelatex
\documentclass[11pt,a4paper]{article}

\usepackage[UTF8]{ctex}
\usepackage[margin=2.2cm]{geometry}
\usepackage{amsmath, amssymb, amsthm}
\usepackage{listings}
\usepackage{xcolor}
\usepackage{tikz}
\usepackage{booktabs}
\usepackage{multirow}
\usepackage{hyperref}
\usetikzlibrary{positioning, arrows.meta, shapes.geometric, fit, calc, decorations.pathreplacing}

\hypersetup{
  colorlinks=true,
  linkcolor=blue!60!black,
  urlcolor=teal!70!black,
}

\lstdefinelanguage{Rust}{
  morekeywords={fn,let,mut,pub,struct,impl,use,mod,self,Self,Option,Some,None,return,if,else,for,while,loop,match,break,continue,as,where,trait,type,enum,move,ref,const,static,unsafe,extern,crate,super,dyn,async,await,box,in},
  sensitive=true,
  morecomment=[l]{//},
  morecomment=[s]{/*}{*/},
  morestring=[b]",
}

\lstset{
  language=Rust,
  basicstyle=\ttfamily\footnotesize,
  keywordstyle=\color{blue!60!black}\bfseries,
  commentstyle=\color{green!40!black}\itshape,
  stringstyle=\color{red!60!black},
  numbers=left,
  numberstyle=\tiny\color{gray},
  numbersep=8pt,
  frame=single,
  framerule=0.4pt,
  rulecolor=\color{gray!40},
  breaklines=true,
  breakatwhitespace=true,
  showstringspaces=false,
  tabsize=4,
  columns=fullflexible,
  keepspaces=true,
}

\title{\textbf{halfedge\,: 属性系统设计文档} \\ \large \texttt{src/property.rs}}
\author{halfedge 库}
\date{2026-07-04}

\begin{document}
\maketitle
\tableofcontents

\section{模块定位}

\texttt{property.rs} 实现 OpenMesh 风格的属性系统，允许用户为顶点/半边/面
附加\textbf{任意类型}的自定义数据。设计目标：

\begin{itemize}
  \item 不依赖 trait 泛型，使用 \texttt{Any + TypeId} 实现类型擦除；
  \item 同一容器内可并存多种类型（\texttt{f64} / \texttt{String} / 自定义 newtype 等）；
  \item 与 \texttt{MeshStorage} \textbf{解耦}，作为独立外部结构体
        \texttt{MeshProperties} 使用；
  \item 删除网格元素时通过包装函数同步清理属性。
\end{itemize}

\section{核心数据结构}

\subsection{PropertyStore<T>：单类型属性存储}

\begin{lstlisting}
pub struct PropertyStore<T> {
    data: HashMap<usize, T>,  // key = slotmap key 的 idx 部分
}
\end{lstlisting}

\textbf{键的选取}：使用 slotmap key 的 \texttt{idx}（槽位索引）部分作为 \texttt{usize} 键，
而非完整 key（包含 version）。原因：

\begin{itemize}
  \item 用户需求明确指定 \texttt{HashMap<usize, T>}；
  \item 同一元素的 idx 在其生命周期内不变；
  \item 便于 \texttt{ErasedProperty::remove\_by\_index} 在不知类型时批量删除。
\end{itemize}

\textbf{权衡}：idx 在 slot 复用时会被新元素继承，因此删除元素后\textbf{必须}
同步清理属性（见 \S\ref{sec:remove-callback}），否则属性会错误关联到新元素。

\subsection{idx 提取辅助函数}

slotmap 1.1.1 的 \texttt{KeyData} 字段为私有，公开 API 仅 \texttt{as\_ffi()}。
通过 \texttt{Key} trait 的 \texttt{data()} 方法取出 \texttt{KeyData}，
再调 \texttt{as\_ffi()} 取 64 位编码，取低 32 位即 idx：

\begin{lstlisting}
#[inline]
fn idx_of<K: Key>(id: K) -> usize {
    let ffi = id.data().as_ffi();
    (ffi & 0xFFFF_FFFF) as usize
}
\end{lstlisting}

\texttt{as\_ffi()} 编码格式（slotmap 源码）：
\[
\texttt{as\_ffi()} = (\underbrace{\texttt{version}}_{\text{高 32 位}} \ll 32)\;\big|\;\underbrace{\texttt{idx}}_{\text{低 32 位}}
\]

\subsection{ErasedProperty trait（私有）}

为支持「不知类型时批量删除」，扩展标准库 \texttt{Any}：

\begin{lstlisting}
trait ErasedProperty {
    fn as_any(&self) -> &dyn Any;
    fn as_any_mut(&mut self) -> &mut dyn Any;
    fn remove_by_index(&mut self, idx: usize);
    fn clear_all(&mut self);
}
\end{lstlisting}

\texttt{as\_any} 与 \texttt{as\_any\_mut} 用于 downcast 回具体
\texttt{PropertyStore<T>}；\texttt{remove\_by\_index} 让容器在
不知 \texttt{T} 的情况下批量删除某 ID 的所有属性。

\subsection{PropertyHandle<T>：类型化句柄}

\begin{lstlisting}
pub struct PropertyHandle<T> {
    _marker: PhantomData<T>,
}
\end{lstlisting}

零大小类型，仅用于编译期类型推断，是 \texttt{Copy} 的。属性按 \texttt{TypeId} 索引，
句柄本身不携带数据，仅让 \texttt{get/set/remove} 方法无需显式标注泛型参数：

\begin{lstlisting}
let h: PropertyHandle<f64> = props.add_vertex_prop();
props.set_vertex_prop(h, v, 1.5);  // T 由 h 推断为 f64
let w = *props.get_vertex_prop(h, v).unwrap();
\end{lstlisting}

\subsection{MeshProperties：属性容器}

\begin{lstlisting}
pub struct MeshProperties {
    vertex_props:   HashMap<TypeId, Box<dyn ErasedProperty>>,
    halfedge_props: HashMap<TypeId, Box<dyn ErasedProperty>>,
    face_props:     HashMap<TypeId, Box<dyn ErasedProperty>>,
}
\end{lstlisting}

三类属性分别存储在独立 \texttt{HashMap} 中，按 \texttt{TypeId::of::<T>()} 区分类型。
每个类型 \texttt{T: 'static} 在每类中只能注册一个属性。

\section{CRUD 方法}

每类元素（vertex / halfedge / face）提供 5 个方法：\texttt{add} / \texttt{get} /
\texttt{get\_mut} / \texttt{set} / \texttt{remove}，加上 \texttt{remove\_*\_all\_props}
批量删除和 \texttt{clear\_*\_props} 全清。以 vertex 为例：

\begin{center}
\renewcommand{\arraystretch}{1.18}
\begin{tabular}{@{}lll@{}}
  \toprule
  \textbf{方法} & \textbf{签名} & \textbf{语义} \\
  \midrule
  \texttt{add\_vertex\_prop::<T>} & \texttt{() -> PropertyHandle<T>} & 注册类型 T \\
  \texttt{get\_vertex\_prop::<T>} & \texttt{(h, id) -> Option<\&T>} & 读取 \\
  \texttt{get\_vertex\_prop\_mut::<T>} & \texttt{(h, id) -> Option<\&mut T>} & 可变读 \\
  \texttt{set\_vertex\_prop::<T>} & \texttt{(h, id, val)} & 写入（自动注册） \\
  \texttt{remove\_vertex\_prop::<T>} & \texttt{(h, id) -> Option<T>} & 删除单类型 \\
  \texttt{remove\_vertex\_all\_props} & \texttt{(id)} & 删除所有类型 \\
  \texttt{has\_vertex\_prop::<T>} & \texttt{() -> bool} & 是否注册 \\
  \texttt{clear\_vertex\_props} & \texttt{()} & 清空但保留类型槽 \\
  \bottomrule
\end{tabular}
\end{center}

\textbf{自动注册}：\texttt{set\_vertex\_prop} 在属性未注册时自动创建
\texttt{PropertyStore<T>} 并插入，简化用户代码。

\section{类型擦除流程}

\subsection{set 写入流程}

\begin{center}
\resizebox{\textwidth}{!}{%
\begin{tikzpicture}[
  node distance=0.45cm,
  proc/.style={draw, rounded corners=2pt, minimum width=11cm, minimum height=0.7cm, align=center, font=\small},
  dec/.style={diamond, draw, aspect=2.4, fill=orange!12, inner sep=1pt, font=\small, align=center, minimum width=4.5cm},
  op/.style={-Stealth, thick},
  startend/.style={draw, rounded corners=10pt, minimum width=2.6cm, minimum height=0.7cm, font=\small, fill=gray!12}
]
  \node[startend] (s) {调用 \texttt{set\_vertex\_prop::<T>(h, id, val)}};
  \node[proc, below=of s, fill=blue!8] (p1) {计算 \texttt{type\_id = TypeId::of::<T>()}};
  \node[dec, below=of p1] (d1) {\texttt{vertex\_props} 是否含 \texttt{type\_id}？};
  \node[proc, below=1.0cm of d1, fill=blue!8] (p2) {downcast\_mut 为 \texttt{PropertyStore<T>}，调用 \texttt{set(idx, val)}};
  \node[proc, right=2.0cm of d1, fill=green!12] (p3) {创建 \texttt{PropertyStore::<T>::new()}，\texttt{set(idx, val)}，\texttt{insert(type\_id, store)}};
  \node[startend, below=of p2, fill=green!12] (ok) {返回};
  \draw[op] (s) -- (p1);
  \draw[op] (p1) -- (d1);
  \draw[op] (d1) -- node[right, font=\footnotesize] {是} (p2);
  \draw[op] (d1) -- node[above, near start, font=\footnotesize] {否} (p3);
  \draw[op] (p2) -- (ok);
  \draw[op] (p3.south) |- (ok.east);
\end{tikzpicture}%
}
\end{center}

\subsection{get 读取流程}

\begin{lstlisting}
pub fn get_vertex_prop<T: 'static>(
    &self,
    _handle: PropertyHandle<T>,
    id: VertexId,
) -> Option<&T> {
    let entry = self.vertex_props.get(&TypeId::of::<T>())?;
    let store = entry.as_any().downcast_ref::<PropertyStore<T>>()?;
    store.get(idx_of(id))
}
\end{lstlisting}

三步：(1) 按 \texttt{TypeId} 查表；(2) downcast 回具体 \texttt{PropertyStore<T>}；
(3) 按 idx 查值。任一环节失败返回 \texttt{None}（属性未注册 / 类型不匹配 / 该 ID 未设值）。

\section{删除回调}
\label{sec:remove-callback}

\texttt{MeshProperties} 与 \texttt{MeshStorage} 解耦，删除网格元素时不会自动清理
对应属性。提供三个包装函数同步清理：

\begin{lstlisting}
pub fn remove_vertex_with_props(
    mesh: &mut MeshStorage,
    props: &mut MeshProperties,
    id: VertexId,
) -> Option<Vertex> {
    let result = mesh.remove_vertex(id);
    props.remove_vertex_all_props(id);  // 遍历所有类型，按 idx 删除
    result
}
\end{lstlisting}

\texttt{remove\_vertex\_all\_props} 通过 \texttt{ErasedProperty::remove\_by\_index}
在不知具体类型的情况下批量删除：

\begin{lstlisting}
pub fn remove_vertex_all_props(&mut self, id: VertexId) {
    let idx = idx_of(id);
    for store in self.vertex_props.values_mut() {
        store.remove_by_index(idx);
    }
}
\end{lstlisting}

\section{应用：加权拉普拉斯平滑}

利用顶点属性 \texttt{f64} 存储权重，执行\textbf{加权}拉普拉斯平滑：

\subsection{公式}

等权拉普拉斯（\texttt{geometry.rs} 中已实现）：
\[
\vec{p}_v' = \frac{1}{|N(v)|} \sum_{u \in N(v)} \vec{p}_u
\]

加权拉普拉斯（用属性权重 $w_u$）：
\[
\vec{p}_v' = \frac{\displaystyle\sum_{u \in N(v)} w_u \cdot \vec{p}_u}{\displaystyle\sum_{u \in N(v)} w_u}
\]

权重 $w_u$ 大的邻居对平滑位置贡献更大，平滑位置偏向高权重邻居。

\subsection{实现要点}

\begin{lstlisting}
let w: PropertyHandle<f64> = props.add_vertex_prop();
for v in mesh.vertex_ids() {
    props.set_vertex_prop(w, v, 1.0);
}
// 把某个邻居权重调大
props.set_vertex_prop(w, neighbor, 10.0);

// 加权平均
let mut sum_w = 0.0;
let mut sum_p = [0.0; 3];
for u in VertexAdjacentVerts::new(&mesh, v) {
    let w_u = *props.get_vertex_prop(w, u).unwrap_or(&1.0);
    let p = mesh.get_vertex(u).unwrap().position;
    for k in 0..3 { sum_p[k] += w_u * p[k]; }
    sum_w += w_u;
}
let weighted_avg = [sum_p[0] / sum_w, sum_p[1] / sum_w, sum_p[2] / sum_w];
\end{lstlisting}

\section{复杂度分析}

\begin{center}
\renewcommand{\arraystretch}{1.20}
\begin{tabular}{@{}lll@{}}
  \toprule
  \textbf{操作} & \textbf{时间复杂度} & \textbf{空间复杂度} \\
  \midrule
  \texttt{add\_*\_prop}（注册类型）       & $O(1)$ & $O(1)$（空 store） \\
  \texttt{get\_*\_prop}（读取）            & $O(1)$ & $O(1)$ \\
  \texttt{set\_*\_prop}（写入）            & $O(1)$ 均摊 & $O(1)$ per entry \\
  \texttt{remove\_*\_prop}（删单类型）     & $O(1)$ & $O(1)$ \\
  \texttt{remove\_*\_all\_props}（删全类型） & $O(K)$，$K$ = 已注册类型数 & $O(1)$ \\
  \texttt{clear\_*\_props}（清空）         & $O(K + N)$，$N$ = 总条目数 & $O(1)$ \\
  \bottomrule
\end{tabular}
\end{center}

\section{单元测试覆盖}

\texttt{src/property.rs} 末尾共 \textbf{16} 个测试，分四组：

\begin{center}
\scriptsize
\renewcommand{\arraystretch}{1.18}
\begin{tabular}{@{}ll@{}}
  \toprule
  \textbf{测试名} & \textbf{覆盖点} \\
  \midrule
  \multicolumn{2}{l}{\textit{PropertyStore 基础}} \\
  \texttt{property\_store\_basic\_crud}       & set / get / get\_mut / remove / contains / len \\
  \texttt{property\_store\_iter}              & 迭代 \texttt{(idx, \&T)} 对 \\
  \texttt{property\_store\_clear}             & clear 后 is\_empty \\
  \midrule
  \multicolumn{2}{l}{\textit{类型擦除}} \\
  \texttt{add\_and\_get\_vertex\_prop}        & 注册 + 读写多顶点 \\
  \texttt{get\_vertex\_prop\_mut\_works}      & 可变引用就地修改 \\
  \texttt{remove\_vertex\_prop}               & 删除单类型 + 重复删除返回 None \\
  \texttt{remove\_vertex\_all\_props\_clears\_all\_types} & 删除顶点时清理所有类型属性 \\
  \texttt{different\_types\_do\_not\_interfere} & f64 / i32 / String 互不干扰 \\
  \texttt{vertex\_halfedge\_face\_props\_are\_independent} & 三类属性独立存储 \\
  \midrule
  \multicolumn{2}{l}{\textit{自动注册与边缘情况}} \\
  \texttt{get\_unregistered\_prop\_returns\_none} & 未注册类型读取返回 None \\
  \texttt{set\_auto\_registers\_prop}         & 未注册时 set 自动注册 \\
  \texttt{halfedge\_and\_face\_props\_work}   & 半边 / 面属性 CRUD + 批量删 \\
  \texttt{clear\_all\_vertex\_props}          & 清空数据但保留类型槽 \\
  \midrule
  \multicolumn{2}{l}{\textit{句柄与包装函数}} \\
  \texttt{remove\_vertex\_with\_props\_wrapper} & 包装函数同步清理属性 \\
  \texttt{property\_handle\_is\_copy\_and\_debug} & 句柄 Copy + Debug \\
  \texttt{newtype\_wrapper\_for\_multiple\_same\_type\_props} & newtype 区分同底层类型属性 \\
  \bottomrule
\end{tabular}
\end{center}

\texttt{cargo test} 全模块共 \textbf{371} 个用例，全部通过（0 失败 / 0 警告）。

\section{设计取舍}

\begin{itemize}
  \item \textbf{为何用 \texttt{Box<dyn ErasedProperty>} 而非 \texttt{Box<dyn Any>}？} \\
        标准库 \texttt{Any} 不提供 \texttt{remove\_by\_index}，
        无法在不知 \texttt{T} 的情况下批量删除。\texttt{ErasedProperty} trait
        扩展了此能力，使 \texttt{remove\_*\_all\_props} 无需为每个已注册类型
        显式调用 \texttt{remove}。
  \item \textbf{为何属性键用 idx 而非完整 key？} \\
        用户需求明确指定 \texttt{HashMap<usize, T>}。idx 在元素生命周期内不变，
        满足稀疏存储需求。代价是删除元素后必须同步清理（见 \S\ref{sec:remove-callback}）。
        若用完整 \texttt{as\_ffi()} 作键，版本号会防止 slot 复用冲突，
        但偏离用户指定接口且增加 \texttt{ErasedProperty} 复杂度。
  \item \textbf{为何 \texttt{MeshProperties} 是外部结构体而非 \texttt{MeshStorage} 字段？} \\
        解耦让 \texttt{MeshStorage} 保持纯粹「存储 + 句柄有效性」职责，
        不绑定具体属性类型。用户可按需创建 0 个或多个 \texttt{MeshProperties}
        实例（如分别存放几何属性与渲染属性）。
  \item \textbf{为何 \texttt{PropertyHandle<T>} 是零大小类型？} \\
        属性按 \texttt{TypeId} 索引，句柄仅用于编译期类型推断。
        零大小 \texttt{Copy} 类型可自由复制，无运行时开销。
  \item \textbf{每类型只能注册一个属性？} \\
        \texttt{TypeId} 作键意味着同类型只能存一份。需要多个同底层类型属性时，
        用 newtype 包装：\texttt{struct Weight(f64); struct Temperature(f64);}
        生成不同 \texttt{TypeId}，互不冲突。
\end{itemize}

\section{使用示例}

\subsection{基础：注册并读写顶点属性}

\begin{lstlisting}
let mesh = build_icosphere(1);
let mut props = MeshProperties::new();

let w: PropertyHandle<f64> = props.add_vertex_prop();
for v in mesh.vertex_ids() {
    props.set_vertex_prop(w, v, 1.0);
}
let v0 = mesh.vertex_ids().next().unwrap();
assert_eq!(props.get_vertex_prop(w, v0), Some(&1.0));
\end{lstlisting}

\subsection{类型擦除：同容器并存多种类型}

\begin{lstlisting}
let label: PropertyHandle<i32>    = props.add_vertex_prop();
let name:  PropertyHandle<String> = props.add_vertex_prop();
let flag:  PropertyHandle<bool>   = props.add_vertex_prop();
// 三种类型互不干扰，vertex_prop_type_count() == 4（含之前的 f64）
\end{lstlisting}

\subsection{newtype：同底层类型多属性}

\begin{lstlisting}
#[derive(Debug, PartialEq)]
struct Weight(f64);
#[derive(Debug, PartialEq)]
struct Temperature(f64);

let hw: PropertyHandle<Weight>      = props.add_vertex_prop();
let ht: PropertyHandle<Temperature> = props.add_vertex_prop();
props.set_vertex_prop(hw, v, Weight(1.0));
props.set_vertex_prop(ht, v, Temperature(25.5));
\end{lstlisting}

\subsection{半边 / 面属性：缓存边长与面法向}

\begin{lstlisting}
let edge_len: PropertyHandle<f64>  = props.add_halfedge_prop();
let normal:   PropertyHandle<Vec3> = props.add_face_prop();

for he in mesh.halfedge_ids() {
    let h = mesh.get_halfedge(he).unwrap();
    let tip = h.vertex;
    let origin = mesh.get_halfedge(h.twin.unwrap()).unwrap().vertex;
    let l = dist(mesh.get_vertex(tip).unwrap().position,
                 mesh.get_vertex(origin).unwrap().position);
    props.set_halfedge_prop(edge_len, he, l);
}
\end{lstlisting}

\subsection{删除元素同步清理}

\begin{lstlisting}
use halfedge::property::remove_vertex_with_props;

let removed = remove_vertex_with_props(&mut mesh, &mut props, vids[0]);
assert!(removed.is_some());
assert!(!mesh.contains_vertex(vids[0]));
assert_eq!(props.get_vertex_prop(w, vids[0]), None);  // 属性已同步清除
\end{lstlisting}

\end{document}
